\section{Explicit constants and proof of the uniform rate}\label{sec:rateproof}

We prove Theorem~\ref{thm:rate} without a lower bound on the moment order or the hyperbolic parameter. Throughout this appendix,
\[
m=\lfloor k/2\rfloor,\quad K=2m,\quad\epsilon=k-K,\quad
\lambda=\frac23\sqrt{K^3/N},\quad
\delta=\frac{K^2}{N},\quad\eta=\frac{3K^2}{4N}.
\]
The symbol $\delta$ here denotes the scale parameter, not parity.

\subsection{Constants and a finite estimate}

For $C>0$, put $a=3C^2/4$ and define
\begin{equation}\label{eq:rateconstantsB}
\begin{aligned}
B_1(C)&=\frac43(2+3a)\sqrt C+\frac{2144}{27}C^{5/2},\\
B_2(C)&=\frac32(2a+3)C
+\frac{16}{9}(245+45a+45a^2)C,\\
B_3(C)&=3(10a^3+188a^2+457a+141).
\end{aligned}
\end{equation}
One valid choice in \eqref{eq:ratefinite} is
\begin{equation}\label{eq:rateconstantsA}
\begin{aligned}
A_1(C)&=e^{7C^2/4}B_1(C)+2e^{11C^2/4}C^{5/2},\\
A_2(C)&=e^{7C^2/4}B_2(C)+2e^{11C^2/4}C
+e^{C^2/3}\left(\frac C2+\frac{C^3}{3}\right),\\
A_3(C)&=e^{7C^2/4}B_3(C).
\end{aligned}
\end{equation}
In particular, the displayed expression for $A_1(C)$ tends to zero as $C\downarrow0$.

We first prove a bound using the actual order $k$. For $2\le k\le N/4$, set
\[
b=\min(\lambda,\lambda^2),\qquad
r=\frac{k}{2N}+\frac{k^3}{3N^2},\qquad
\tau=\frac{2K^{5/2}}{N^{3/2}},\qquad\zeta=\frac{2K}{N},
\]
and
\begin{equation}\label{eq:rateB}
\begin{aligned}
\mathcal B={}&\frac{(2+3\eta)b+2\eta^2+3\eta}{m}\\
&+\frac{1}{m^2}\bigl[67\lambda^3+(245+45\eta+45\eta^2)\lambda^2\\
&\hspace{27mm}+10\eta^4+188\eta^3+457\eta^2+141\eta\bigr].
\end{aligned}
\end{equation}
The finite estimate is
\begin{equation}\label{eq:ratelocal}
|W_{N,k}-1|
\le e^r-1+D_ke^{-k^2/(4N)+\eta}
\bigl[\mathcal B+(1+\zeta)e^\tau-1\bigr].
\end{equation}

\subsection{Condition the triple count on its parity}

Let $J$ be Poisson with mean $\lambda$, conditioned to have parity $\epsilon$. Independently let $V$ be Poisson with mean $\eta$. Put
\[
T=J-\epsilon,\qquad S=3T/2+2V,
\]
and define
\[
w_\epsilon(s)=\frac{(m-\epsilon)_s}{m^s}
\quad(0\le s\le m-\epsilon),
\]
with value zero at larger nonnegative integers. The variables $T$ and $S$ are nonnegative integers. Restricting \eqref{eq:availabilityeven} and \eqref{eq:availabilityodd} to blocks of sizes two, three and four shows that their contribution to $(Q_k/\mathcal L_k)/H_k$ is exactly
\begin{equation}\label{eq:rateconditioned}
e^\eta\E w_\epsilon(S).
\end{equation}
At odd orders, the $1/\lambda$ factors in the coefficient formula and $H_k$ cancel. Thus \eqref{eq:rateconditioned} is valid even when $\lambda$ is small.

The following bounds hold for both conditioned laws:
\begin{align}
\E T&\le b,&
|\E T^2-\lambda^2|&\le b,&
\E T^2&\le2\lambda^2,\label{eq:ratemomentsone}\\
\E T^3&\le\lambda^3+4\lambda^2,&
\E(T^2-\lambda^2)^2&\le4\lambda^3+8\lambda^2.
&&\label{eq:ratemomentstwo}
\end{align}
For verification, write $v=\lambda\tanh\lambda$ and $q=\lambda\coth\lambda-1$. Differentiating the normalizers $\cosh\lambda$ and $\sinh\lambda$ gives
\[
\begin{array}{c|cc}
&\epsilon=0&\epsilon=1\\ \hline
\E T&v&q\\
\E T^2&\lambda^2+v&\lambda^2-q\\
\E T^3&\lambda^2v+3\lambda^2+v&\lambda^2+\lambda^2q+q\\
\E(T^2-\lambda^2)^2&4\lambda^2v+7\lambda^2+v&3\lambda^2+(4\lambda^2-1)q
\end{array}
\]
Both $v$ and $q$ lie in $[0,b]$. For $q$, use $\coth\lambda\le1+1/\lambda$ and
$\lambda\cosh\lambda\le(1+\lambda^2)\sinh\lambda$. The latter inequality follows by differentiation, since the difference is zero at zero and has derivative $\lambda^2\cosh\lambda+\lambda\sinh\lambda\ge0$. The table now proves \eqref{eq:ratemomentsone} and \eqref{eq:ratemomentstwo}.

\subsection{Compare pair availability with an exponential}

For an integer $s\ge0$, set $A(s)=s(s-1+2\epsilon)/(2m)$. Then
\begin{equation}\label{eq:rateavailability}
0\le e^{-A(s)}-w_\epsilon(s)\le4s^3/m^2.
\end{equation}
If $s\le m/2$, expand the logarithm of the product. Its linear sum is $-A(s)$, and
\[
0\le-\log(1-u)-u\le\frac{u^2}{2(1-u)}
\]
bounds the remaining sum by $s^3/m^2$. Use $1-e^{-z}\le z$ to obtain \eqref{eq:rateavailability}. If $s>m/2$, the difference lies in $[0,1]$ and $4s^3/m^2>1$, including when the product is assigned zero.

Write $A=A(S)$. Since $\delta=9\lambda^2/(8m)$,
\begin{equation}\label{eq:rateAexpand}
A-\delta
=\frac{9(T^2-\lambda^2)}{8m}+\frac{3TV}{m}+\frac{2V^2}{m}
+(2\epsilon-1)\left(\frac{3T}{4m}+\frac Vm\right).
\end{equation}
Independence and the Poisson moments
\[
\E V^2=\eta^2+\eta,\quad
\E V^3=\eta^3+3\eta^2+\eta,\quad
\E V^4=\eta^4+6\eta^3+7\eta^2+\eta
\]
give
\begin{equation}\label{eq:rateAmean}
|\E A-\delta|
\le\frac{(15/8+3\eta)b+2\eta^2+3\eta}{m}.
\end{equation}
Apply $(u_1+\cdots+u_5)^2\le5\sum_i u_i^2$ to the five terms in \eqref{eq:rateAexpand}. Equations~\eqref{eq:ratemomentsone} and \eqref{eq:ratemomentstwo} then imply
\begin{equation}\label{eq:rateAsecond}
\begin{aligned}
\E(A-\delta)^2\le\frac5{m^2}\bigl[&\frac{81}{16}\lambda^3
+(45/4+18\eta+18\eta^2)\lambda^2\\
&+4\eta^4+24\eta^3+29\eta^2+5\eta\bigr].
\end{aligned}
\end{equation}
Also $(u+v)^3\le4(u^3+v^3)$ for nonnegative $u,v$ gives
\begin{equation}\label{eq:rateSthird}
\E S^3\le\frac{27}{2}(\lambda^3+4\lambda^2)
+32(\eta^3+3\eta^2+\eta).
\end{equation}

Because $A,\delta\ge0$ and the second derivative of $e^{-u}$ has absolute value at most one on this interval, Taylor's theorem yields
\[
|\E e^{-A}-e^{-\delta}|
\le|\E A-\delta|+\tfrac12\E(A-\delta)^2.
\]
Combine this with \eqref{eq:rateavailability} and \eqref{eq:rateAmean} through \eqref{eq:rateSthird}. Rounding the positive coefficients upward gives
\begin{equation}\label{eq:ratewerror}
|\E w_\epsilon(S)-e^{-\delta}|\le\mathcal B.
\end{equation}
For example, the coefficient of $\lambda^3/m^2$ before rounding is $54+405/32=2133/32<67$, and that of $\lambda^2/m^2$ independent of $\eta$ is $1953/8<245$.

\subsection{Larger blocks and the denominator}

Let $x=\sqrt{K/N}\le1/2$ and $\lambda_j=N(j-1)x^j/j$. Denote by $T_{\rm e}$ and $T_{\rm o}$ the sums of these weights over even and odd $j\ge5$. Then
\[
T_{\rm e}+T_{\rm o}\le\frac{Nx^5}{1-x}\le\tau,
\qquad
\frac{T_{\rm o}}\lambda\le\frac{3x^2}{2(1-x^2)}\le\zeta.
\]
All coefficients are nonnegative. Removing availability factors bounds the additional contribution after division by $H_k$ by
\[
e^\eta(e^\tau-1)
\quad\text{at even orders},\qquad
e^\eta\bigl[(1+\zeta)e^\tau-1\bigr]
\quad\text{at odd orders}.
\]
Indeed, summing the added odd weights replaces $\cosh\lambda$ or $\sinh\lambda$ by its value at $\lambda+T_{\rm o}$. The logarithmic derivative of $\cosh u$ is at most one, as is that of $\sinh u/u$. The latter gives the extra factor $1+T_{\rm o}/\lambda$ and preserves the small-$\lambda$ normalization.

Together with \eqref{eq:rateconditioned} and \eqref{eq:ratewerror}, this proves
\begin{equation}\label{eq:rateQerror}
\left|\frac{Q_k/\mathcal L_k}{H_k}-e^{\eta-\delta}\right|
\le e^\eta\bigl[\mathcal B+(1+\zeta)e^\tau-1\bigr].
\end{equation}
Since $\eta-\delta=-K^2/(4N)$, the central value after restoring $D_k$ and the correction in \eqref{eq:rateW} is
$D_k\exp(-(k^2+K^2)/(4N))$.
The logarithm remainder inequality gives, for $k\le N/4$,
\[
0\le\log D_k-\frac{k(k-1)}{2N}\le\frac{k^3}{3N^2}.
\]
The difference between $k(k-1)/(2N)$ and $(k^2+K^2)/(4N)$ is $-k/(2N)$ for even $k$ and $-1/(4N)$ for odd $k$. Thus its total logarithmic error has absolute value at most $r$. Multiplying \eqref{eq:rateQerror} by $D_ke^{-k^2/(4N)}$ proves \eqref{eq:ratelocal}.

\subsection{Uniform bounds for every positive endpoint}

Assume $C>0$, $N\ge16C^2$ and $2\le k\le C\sqrt N$. Then $k\le N/4$, $K\le C\sqrt N$, $\eta\le a$ and $\delta\le C^2$. The identities
\[
\frac\lambda m=\frac43\sqrt{K/N},\quad
\frac\eta m=\frac{3K}{2N},\quad
\frac{\lambda^3}{m^2}=\frac{32K^{5/2}}{27N^{3/2}},\quad
\frac{\lambda^2}{m^2}=\frac{16K}{9N},\quad
\frac\eta{m^2}=\frac3N
\]
show, using $b\le\lambda$, that
\[
\mathcal B\le B_1(C)N^{-1/4}+B_2(C)N^{-1/2}+B_3(C)N^{-1}.
\]
Moreover, $\tau\le\delta\le C^2$, so
\[
(1+\zeta)e^\tau-1
\le e^{C^2}(\tau+\zeta)
\le2e^{C^2}\bigl(C^{5/2}N^{-1/4}+CN^{-1/2}\bigr).
\]
The multiplier $D_ke^{-k^2/(4N)+\eta}$ in \eqref{eq:ratelocal} is at most $e^{7C^2/4}$ because $\log D_k\le k^2/N$. Finally,
\[
r\le\left(\frac C2+\frac{C^3}{3}\right)N^{-1/2},\qquad
r=\frac{k^2}{N}\left(\frac1{2k}+\frac{k}{3N}\right)\le\frac{C^2}{3}.
\]
Hence $e^r-1\le e^{C^2/3}(C/2+C^3/3)N^{-1/2}$. Substitution in \eqref{eq:ratelocal} gives precisely \eqref{eq:rateconstantsA} and proves \eqref{eq:ratefinite}. These comparisons use only $C>0$, not $C\ge1$. They also justify the small-endpoint limit of $A_1$ used in Theorem~\ref{thm:firstcorrection}.
